\documentclass[10pt,a4paper]{article}
\usepackage[utf8]{inputenc}
\usepackage{graphicx}
\usepackage[left=2.54cm, right=1.0in, top=1.0in, bottom=1.0in]{geometry}
\usepackage[swedish]{babel}
\usepackage{tikz}
\usetikzlibrary{arrows.meta,shapes,calc,backgrounds}
\usepackage{parskip}
\usepackage{xcolor}
\usepackage[most]{tcolorbox}
\usepackage{mdframed}
\usepackage{tabularx}
\usepackage{colortbl}
\usepackage{enumerate}
\usepackage{enumitem}

\setlength{\extrarowheight}{8pt}
\renewcommand{\arraystretch}{1}

% My colors
\definecolor{mycolor1}{rgb}{1,0,0} %red
\definecolor{lightpurple}{HTML}{c69cd6}
\definecolor{darkpurple}{HTML}{7b5c87}

%circled
\newcommand*\circled[1]{\tikz[baseline=(char.base)]{
		\node[shape=circle,draw,minimum size=1em, inner sep=2pt] (char)
		{#1};}
}

% Purpleleftline
\newmdenv[innerleftmargin=20,topline=false,rightline=false,bottomline=false, linewidth=3pt, linecolor=darkpurple]{purpleleftline}

% Kommentarboxen
\tikzstyle{mybox} = [draw=lightpurple!35, fill=lightpurple!35, very thick,
rectangle, inner sep=15pt, inner ysep=15pt]
\tikzstyle{fancytitle} =[fill=darkpurple, text=white]

% Title format
\newcommand{\pagetitle}[2] {
	\begin{tabularx}{\linewidth}{>{\columncolor{darkpurple}}c >{\columncolor{lightpurple!35}}X}
		\LARGE \textcolor{white}{{#1}} & \LARGE {#2} \\[3pt]
	\end{tabularx}
}

% Instruction environment
\newenvironment{instructionsenv}
{
	\vspace{-4pt}
	\begin{purpleleftline}
		\vspace{2em}
		{\Large\bfseries{Instruktioner}}\par
		\vspace{1em}
	}
	{ 
	\end{purpleleftline}
	\vspace{50pt}
}

% Instruction
\newcommand{\instruction}[2] {
	\textcolor{mycolor1}{\circled{\large {#1}}} 
	{\large #2}\par
}

% Kommentarbox
\newcommand{\kommentarbox}[1]{
	\begin{center}
		\begin{tikzpicture}
		\node [mybox] (box){%
			\begin{minipage}{0.91\textwidth}
			{#1}
			\end{minipage}
		};
		\node[fancytitle, right=15pt] at (box.north west) {\bfseries Kommentar};
		\end{tikzpicture}
	\end{center}	
}

%\begin{bild}{fil}{storlek andel av linewidth}
%\end{bild}
\newenvironment{bild}[2]
{
\begin{center}
	\begin{tikzpicture}
	\node[anchor=south west,inner sep=0] (image) at (0,0) {\includegraphics[width=#2\textwidth]{./img/#1}};
	\begin{scope}[x={(image.south east)},y={(image.north west)}]
}
{		
	\end{scope}
	\end{tikzpicture}
\end{center}
\vspace{35pt}
}

%\siffrastreckcirkel{siffra}{x}{y}{storlek cirkel}{x}{y}
\newcommand{\siffrastreckcirkel}[6]{
	\node[mycolor1, thick] (A) at  ({#2},{#3}) {\Large {#1}};
	\node[circle, draw, mycolor1, thick, minimum size=#4cm] (B) at  ({#5},{#6}) {};
	\draw[mycolor1, thick] (A) -- (B);
}

% \siffrastreckellips{siffra}{x}{y}{längd ellips}{höjd ellips}{x}{y}
\newcommand{\siffrastreckellips}[7]{
	\node[mycolor1, thick] (A) at  ({#2},{#3}) {\LARGE {#1}};
	\node[ellipse, draw, mycolor1,thick,minimum height=#5cm,minimum width=#4cm,draw] (B) at  ({#6},{#7}) {};
\draw[mycolor1, thick] (A) -- (B);
}

% \siffrapil{siffra}{från x}{från y}{till x}{till y}
\newcommand{\siffrapil}[5]{
	\node[mycolor1, thick] (A) at  (#2,#3) {\LARGE #1};
	\node (B) at  (#4,#5) {};
\draw[-Stealth, mycolor1, line width=1pt] (A)--(B);
}

\begin{document}

%%%%%%%%%%%%%%%%%%%%%%%%%%%%%%%%%%%%%%%%%%%%%%%%%%%%%%%%%%%%%%
%%% Titelsidan
%%%%%%%%%%%%%%%%%%%%%%%%%%%%%%%%%%%%%%%%%%%%%%%%%%%%%%%%%%%%%%
\pagenumbering{gobble}
\newgeometry{left=0cm,bottom=0cm,right=0cm,top=0cm}
\begin{tikzpicture}

\fill[lightpurple!35] (0,0) rectangle (\paperwidth,\paperheight);

\fill[white] (0,0) rectangle (0.65\paperwidth,\paperheight);

\fill[darkpurple] (0,0.45\paperheight) rectangle (\paperwidth,0.65\paperheight);

\node[align=left,white] at (0.5\paperwidth,0.57\paperheight) 
{\bfseries\Huge Exempel på användande av \\ \bfseries\Huge EMG-utrustningen på HSC};

\node[align=center,white] at (0.5\paperwidth,0.505\paperheight) 
{ \huge En praktisk guide};

\node[align=center,black] at (0.825\paperwidth,0.1\paperheight) 
{\Large Peter Kvillegård \\[3em] Version 1.0 \\ 2020-01-23};

\end{tikzpicture}
\restoregeometry
\newpage

%%%%%%%%%%%%%%%%%%%%%%%%%%%%%%%%%%%%%%%%%%%%%%%%%%%%%%%%%%%%%%
%%% Bakgrund
%%%%%%%%%%%%%%%%%%%%%%%%%%%%%%%%%%%%%%%%%%%%%%%%%%%%%%%%%%%%%%

\pagetitle{Om denna guide}{}\par
\vspace{1em}
Jag och Jonatan Malmström använde under hösten 2019 EMG-utrustningen på HSC i Lund för ett experiment i samband med vår C-uppsats. Vi lärde oss att använda program och apparater av andra elever och av vår handledare. I denna användarguide har jag skrivit ner hur vi använde utrustningen, i förhoppning om att det kan vara användbart för andra.

Utrustningen som användes och beskrivs här är MegaWin 3.1-b10 mjukvara med Mega ME6000-T16 biomonitor.  

Vårt arbete fokuserade på att få fram medelamplitud under en definierad tid i olika muskler samtidigt, och användarguiden är därför inriktad för det ändamålet. Detta är inte är en allomfattande manual, utan ett exempel på hur utrustningen kan användas. Det är mycket möjligt att det finns bättre och smidigare sätt, och det finns många andra ändamål som inte kommer att beskrivas här.

Guiden kommer att ta dig steg för steg från att logga in på datorn till resultat som kan användas i er studie. För mer information om EMG-experiment kan jag rekommendera följande böcker:
\begin{itemize}
	\item Konrad, P. The ABC of EMG version 1.4. 2006.
	\item Criswell, E. Cram's Introduction to Surface Electromyography. 2010.
\end{itemize}

Denna guide är tillgänglig under Creative Commons Zero (CC0), vilket lite förenklat innebär att vem som helst har rätt att göra vad som helst med innehållet. Du får använda, modifiera, och distribuera materialet utan begränsningar och utan krav på att attribuera mig. Källkoden är skriven i \LaTeX{} och finns att hitta tillsammans med bilderna på \verb|http://www.github.com/mouboo/lund_emg_userguide|. Om du vill korrigera fel och/eller bygga vidare på denna guide är du mycket välkommen att göra så. Behöver du hjälp med detta, kontakta mig på \verb|peterkvillegård@gmail.com| för mer information.

\newpage

%%%%%%%%%%%%%%%%%%%%%%%%%%%%%%%%%%%%%%%%%%%%%%%%%%%%%%%%%%%%%%
%%% Innehåll
%%%%%%%%%%%%%%%%%%%%%%%%%%%%%%%%%%%%%%%%%%%%%%%%%%%%%%%%%%%%%%

\pagetitle{Innehåll}{}\par
\vspace{3em}
{\Large

\begin{enumerate}[label={\Roman*},align=left]
\item Logga in \dotfill 1
\item Skapa protokoll \dotfill 2
\item Skapa person \dotfill 13
\item Anslut biomonitor \dotfill 16
\item Kontrollera signalen \dotfill 17
\item Utför mätning \dotfill 21
\item Bearbeta signalen \dotfill 25
\end{enumerate}
}

\newpage
\pagenumbering{arabic}

%%%%%%%%%%%%%%%%%%%%%%%%%%%%%%%%%%%%%%%%%%%%%%%%%%%%%%%%%%%%%%
%%% Logga in | Steg 1/1
%%%%%%%%%%%%%%%%%%%%%%%%%%%%%%%%%%%%%%%%%%%%%%%%%%%%%%%%%%%%%%

%%%%% Title
\pagetitle{I. Logga in}{Steg 1/1}

%%%%% Instruktioner
\begin{instructionsenv}
	\instruction{1}{Kryssa i ``Workstation only''.}
	\instruction{2}{Fyll i Username: \texttt{emg}, Password: \texttt{emgemg}.}
	\instruction{3}{Klicka på ``OK''.}
\end{instructionsenv}

%%%%% Bilden
\begin{bild}{login}{0.5}
	%\siffrastreckcirkel{siffra}{x}{y}{storlek cirkel}{x}{y}
	\siffrastreckcirkel{1}{0.2}{0.1}{0.5}{0.03}{0.22}	

	%\siffrastreckcirkel{siffra}{x}{y}{storlek cirkel}{x}{y}
	\siffrastreckcirkel{2}{0.5}{0.6}{1.2}{0.22}{0.43}

	% \siffrastreckellips{siffra}{x}{y}{längd ellips}{höjd ellips}{x}{y}
\siffrastreckellips{3}{0.6}{0.25}{1.5}{0.5}{0.39}{0.1}
\end{bild}

\newpage

%%%%%%%%%%%%%%%%%%%%%%%%%%%%%%%%%%%%%%%%%%%%%%%%%%%%%%%%%%%%%%
%%% II. Skapa protokoll | Steg 1/11
%%%%%%%%%%%%%%%%%%%%%%%%%%%%%%%%%%%%%%%%%%%%%%%%%%%%%%%%%%%%%%

%%%%% Title
\pagetitle{II. Skapa protokoll}{Steg 1/11}

%%%%% Instruktioner
\begin{instructionsenv}
	\instruction{1}{Klicka på ``Protocol''.}
\end{instructionsenv}

%%%%% Bilden
\begin{bild}{3}{0.91}
	%\siffrastreckcirkel{siffra}{x}{y}{storlek cirkel}{x}{y}
	\siffrastreckcirkel{1}{0.25}{0.65}{2}{0.32}{0.85}
\end{bild}

%%%%% Kommentarboxen
\kommentarbox{
	Du behöver bara skapa protokoll en gång, detta kan sedan användas för alla kommande mätningar.
}
\newpage

%%%%%%%%%%%%%%%%%%%%%%%%%%%%%%%%%%%%%%%%%%%%%%%%%%%%%%%%%%%%%%
%%% II. Skapa protokoll | Steg 2/11
%%%%%%%%%%%%%%%%%%%%%%%%%%%%%%%%%%%%%%%%%%%%%%%%%%%%%%%%%%%%%%

%%%%% Title
\pagetitle{II. Skapa protokoll}{Steg 2/11}

%%%%% Instruktioner
\begin{instructionsenv}
	\instruction{1}{Klicka på ``New...''}
\end{instructionsenv}

%%%%% Bilden
\begin{bild}{7}{0.8}
	\siffrastreckellips{1}{0.75}{0.75}{1.5}{0.6}{0.92}{0.695}
\end{bild}

\newpage

%%%%%%%%%%%%%%%%%%%%%%%%%%%%%%%%%%%%%%%%%%%%%%%%%%%%%%%%%%%%%%
%%% II. Skapa protokoll | Steg 3/11
%%%%%%%%%%%%%%%%%%%%%%%%%%%%%%%%%%%%%%%%%%%%%%%%%%%%%%%%%%%%%%

%%%%% Title
\pagetitle{II. Skapa protokoll}{Steg 3/11}

%%%%% Instruktioner
\begin{instructionsenv}
	\instruction{1}{Fyll i protokollets namn.}
	\instruction{2}{Välj ``Advanced protocol''.}
	\instruction{3}{Klicka på ``Next''.}
\end{instructionsenv}

%%%%% Bilden
\begin{bild}{8}{0.8}
	%\siffrastreckcirkel{siffra}{x}{y}{storlek cirkel}{x}{y}	
	\siffrastreckcirkel{1}{0.55}{0.7}{1}{0.25}{0.8}

	% \siffrastreckellips{siffra}{x}{y}{längd ellips}{höjd ellips}{x}{y}
	\siffrastreckellips{2}{0.5}{0.55}{3.5}{0.8}{0.15}{0.565}

	% \siffrastreckellips{siffra}{x}{y}{längd ellips}{höjd ellips}{x}{y}
\siffrastreckellips{3}{0.5}{0.15}{2.5}{0.7}{0.495}{0.037}
\end{bild}

%%%%% Kommentarboxen
\kommentarbox{
	I denna guide har jag valt att kalla protokollet ``ledare'', döp gärna ert protokoll till något unikt som ni kommer ihåg, exempelvis försöksledarnas initialer.
}
\newpage

%%%%%%%%%%%%%%%%%%%%%%%%%%%%%%%%%%%%%%%%%%%%%%%%%%%%%%%%%%%%%%
%%% II. Skapa protokoll | Steg 4/11
%%%%%%%%%%%%%%%%%%%%%%%%%%%%%%%%%%%%%%%%%%%%%%%%%%%%%%%%%%%%%%

%%%%% Title
\pagetitle{II. Skapa protokoll}{Steg 4/11}

%%%%% Instruktioner
\begin{instructionsenv}
	\instruction{1}{Välj ``ME6000-T16''.}
	\instruction{2}{Klicka på ``Next''.}
\end{instructionsenv}

%%%%% Bilden
\begin{bild}{9}{0.8}
	% \siffrastreckellips{siffra}{x}{y}{längd ellips}{höjd ellips}{x}{y}
	\siffrastreckellips{1}{0.55}{0.8}{2}{0.4}{0.2}{0.805}

	% \siffrastreckellips{siffra}{x}{y}{längd ellips}{höjd ellips}{x}{y}
	\siffrastreckellips{2}{0.5}{0.15}{2.5}{0.7}{0.495}{0.037}
\end{bild}

\newpage

%%%%%%%%%%%%%%%%%%%%%%%%%%%%%%%%%%%%%%%%%%%%%%%%%%%%%%%%%%%%%%
%%% II. Skapa protokoll | Steg 5/11
%%%%%%%%%%%%%%%%%%%%%%%%%%%%%%%%%%%%%%%%%%%%%%%%%%%%%%%%%%%%%%

%%%%% Title
\pagetitle{II. Skapa protokoll}{Steg 5/11}

%%%%% Instruktioner
\begin{instructionsenv}
	\instruction{1}{Välj ``General body'', eller annat alternativ.}
	\instruction{2}{Klicka på ``Next''.}
\end{instructionsenv}

%%%%% Bilden
\begin{bild}{10}{0.8}
	% \siffrastreckellips{siffra}{x}{y}{längd ellips}{höjd ellips}{x}{y}
	\siffrastreckellips{1}{0.5}{0.72}{4.5}{1.5}{0.22}{0.73}

	% \siffrastreckellips{siffra}{x}{y}{längd ellips}{höjd ellips}{x}{y}
	\siffrastreckellips{2}{0.5}{0.15}{2.5}{0.7}{0.495}{0.037}
\end{bild}

%%%%% Kommentarboxen
\kommentarbox{
	Ni kan även välja att begränsa vilka muskler som ska visas genom att välja ett annat alternativ. Vi hittade inget sätt att skapa ett protokoll som använder muskler från olika kategorier här, därför rekommenderar vi att välja ``General body'', eller den kategori som passar bäst för ert experiment, och vid behov lägga till muskler som saknas i kommande steg.
}

\newpage

%%%%%%%%%%%%%%%%%%%%%%%%%%%%%%%%%%%%%%%%%%%%%%%%%%%%%%%%%%%%%%
%%% II. Skapa protokoll | Steg 6/11
%%%%%%%%%%%%%%%%%%%%%%%%%%%%%%%%%%%%%%%%%%%%%%%%%%%%%%%%%%%%%%

%%%%% Title
\pagetitle{II. Skapa protokoll}{Steg 6/11}

%%%%% Instruktioner
\begin{instructionsenv}
	\instruction{1}{Markera första kanalen.}
	\instruction{2}{Klicka på ``Select source...''}
\end{instructionsenv}

%%%%% Bilden
\begin{bild}{11}{0.8}
	\siffrastreckellips{1}{0.2}{0.79}{5}{0.7}{0.58}{0.79}
	\siffrastreckellips{2}{0.2}{0.67}{3}{1}{0.2}{0.53}
\end{bild}

%%%%% Kommentarboxen
\kommentarbox{
	Kanalerna motsvarar samma ordning som sladdarna är märkta. ME6000-T16 har 4 ingångar som vardera kan ha 4 kanaler. 
}

\newpage

%%%%%%%%%%%%%%%%%%%%%%%%%%%%%%%%%%%%%%%%%%%%%%%%%%%%%%%%%%%%%%
%%% II. Skapa protokoll | Steg 7/11
%%%%%%%%%%%%%%%%%%%%%%%%%%%%%%%%%%%%%%%%%%%%%%%%%%%%%%%%%%%%%%

\pagetitle{II. Skapa protokoll}{Steg 7/11}

\begin{instructionsenv}
	\instruction{1}{Välj om muskeln är på höger eller vänster sida.}
	\instruction{2}{Markera rätt muskel.}
	\instruction{3}{Klicka på ``OK''.}
\end{instructionsenv}

\vspace{-20pt}

\begin{bild}{12}{0.8}
	\siffrastreckellips{1}{0.15}{1.1}{3.5}{0.7}{0.125}{0.933}
	\siffrastreckellips{2}{-0.07}{0.603}{4}{0.5}{0.15}{0.603}
	\node at (1.07,0.5) {\textcolor{white}{.}};
	\siffrastreckellips{3}{0.675}{0.18}{2.6}{0.8}{0.675}{0.05}

\end{bild}

\kommentarbox{
	Muskeln som ni väljer bestämmer vilket namn som kommer att stå för denna kanal under mätningar och resultat. Om ni inte hittar rätt muskel kan ni välja ``Add Muscle'' och skriva in ett eget namn. Detta är alltså enbart för att ni ska ha koll på vilken kanal som mäter vad i framtida steg.  
}
\newpage

%%%%%%%%%%%%%%%%%%%%%%%%%%%%%%%%%%%%%%%%%%%%%%%%%%%%%%%%%%%%%%
%%% II. Skapa protokoll | Steg 8/11
%%%%%%%%%%%%%%%%%%%%%%%%%%%%%%%%%%%%%%%%%%%%%%%%%%%%%%%%%%%%%%

\pagetitle{II. Skapa protokoll}{Steg 8/11}

\begin{instructionsenv}
	\instruction{1}{Om du vill lägga till fler muskler, upprepa steg 6 och 7 med nästa fria kanal.}
	\instruction{2}{När du har lagt till alla muskler, klicka på ``Next''.}
\end{instructionsenv}

\begin{bild}{13}{0.8}
	\siffrastreckellips{1}{0.2}{0.75}{5}{0.75}{0.58}{0.75}
	\siffrastreckellips{2}{0.37}{0.13}{2.6}{0.8}{0.495}{0.037}
\end{bild}

\newpage

%%%%%%%%%%%%%%%%%%%%%%%%%%%%%%%%%%%%%%%%%%%%%%%%%%%%%%%%%%%%%%
%%% II. Skapa protokoll | Steg 9/11
%%%%%%%%%%%%%%%%%%%%%%%%%%%%%%%%%%%%%%%%%%%%%%%%%%%%%%%%%%%%%%

\pagetitle{II. Skapa protokoll}{Steg 9/11}

\begin{instructionsenv}
	\instruction{1}{Klicka på ``Raw Free''.}
	\instruction{2}{Klicka på ``Pick to protocol''.}
\end{instructionsenv}

\begin{bild}{14}{0.8}
	\siffrastreckellips{1}{0.4}{0.8}{1.5}{1.5}{0.325}{0.645}
	\siffrastreckellips{2}{0.1}{0.5}{3.7}{0.8}{0.498}{0.49}
\end{bild}

\newpage

%%%%%%%%%%%%%%%%%%%%%%%%%%%%%%%%%%%%%%%%%%%%%%%%%%%%%%%%%%%%%%
%%% II. Skapa protokoll | Steg 10/11
%%%%%%%%%%%%%%%%%%%%%%%%%%%%%%%%%%%%%%%%%%%%%%%%%%%%%%%%%%%%%%

\pagetitle{II. Skapa protokoll}{Steg 10/11}

\begin{instructionsenv}
	\instruction{1}{Avmarkera ``Stop when time is up''.}
	\instruction{2}{Klicka på ``OK''.}
\end{instructionsenv}

\begin{bild}{15}{0.8}
	\siffrastreckellips{1}{0.46}{0.32}{0.5}{0.5}{0.545}{0.25}
	\siffrastreckellips{2}{0.78}{0.13}{1.5}{0.8}{0.92}{0.035}
\end{bild}

\kommentarbox{
	Här finns möjlighet att ändra samplingsfrekvensen. En tumregel är att ha högre frekvens än dubbla övre gränsen för bandpassfrekvenserna. 1000 Hz är ett rimligt värde för ett bandpassfilter på 15--400 Hz.\par
	``Time scale'' avgör hur mycket tid som kommer visas i fönstret vid en mätning samtidigt. Med ``Stop when time is up'' avmarkerad så kommer mätningen att fortsätta även efter denna tid, tills den avbryts av operatören. 
}
\newpage

%%%%%%%%%%%%%%%%%%%%%%%%%%%%%%%%%%%%%%%%%%%%%%%%%%%%%%%%%%%%%%
%%% II. Skapa protokoll | Steg 11/11
%%%%%%%%%%%%%%%%%%%%%%%%%%%%%%%%%%%%%%%%%%%%%%%%%%%%%%%%%%%%%%

\pagetitle{II. Skapa protokoll}{Steg 11/11}

\begin{instructionsenv}
	\instruction{1}{Klicka på ``Finish''.}
\end{instructionsenv}

\begin{bild}{16}{0.8}
	\siffrastreckellips{1}{0.7}{0.15}{2}{0.8}{0.7}{0.035}
\end{bild}

\kommentarbox{
	Du har nu skapat ett protokoll.
}
\newpage

%%%%%%%%%%%%%%%%%%%%%%%%%%%%%%%%%%%%%%%%%%%%%%%%%%%%%%%%%%%%%%
%%% III. Skapa person | Steg 1/3
%%%%%%%%%%%%%%%%%%%%%%%%%%%%%%%%%%%%%%%%%%%%%%%%%%%%%%%%%%%%%%

\pagetitle{III. Skapa person}{Steg 1/3}

\begin{instructionsenv}
	\instruction{1}{Klicka på ``Person''.}
\end{instructionsenv}

\begin{bild}{3}{0.91}
	\siffrastreckellips{1}{0.14}{0.6}{2}{2}{0.14}{0.85}
\end{bild}

\kommentarbox{
	Du behöver skapa en person för varje testperson. Det är ett snabbt steg som kan göras vid testtillfället. 
}
\newpage

%%%%%%%%%%%%%%%%%%%%%%%%%%%%%%%%%%%%%%%%%%%%%%%%%%%%%%%%%%%%%%
%%% III. Skapa person | Steg 2/3
%%%%%%%%%%%%%%%%%%%%%%%%%%%%%%%%%%%%%%%%%%%%%%%%%%%%%%%%%%%%%%

\pagetitle{III. Skapa person}{Steg 2/3}

\begin{instructionsenv}
	\instruction{1}{Klicka på ``New person''.}
\end{instructionsenv}

\begin{bild}{4}{0.8}
	\siffrastreckellips{1}{0.7}{0.6}{2.2}{0.8}{0.89}{0.57}
\end{bild}

\newpage
%%%%%%%%%%%%%%%%%%%%%%%%%%%%%%%%%%%%%%%%%%%%%%%%%%%%%%%%%%%%%%
%%% III. Skapa person | Steg 3/3
%%%%%%%%%%%%%%%%%%%%%%%%%%%%%%%%%%%%%%%%%%%%%%%%%%%%%%%%%%%%%%

\pagetitle{III. Skapa person}{Steg 3/3}

\begin{instructionsenv}
	\instruction{1}{Fyll i ``Last name''. Resten behöver inte fyllas i.}
	\instruction{2}{Klicka på ``Save''.}
\end{instructionsenv}

\begin{bild}{5}{0.8}
	\siffrastreckellips{1}{0.05}{0.775}{3.3}{0.5}{0.245}{0.775}
	\siffrastreckellips{2}{0.7}{0.2}{2.2}{0.8}{0.885}{0.19}
\end{bild}

\kommentarbox{
	Med tanke på testpersonernas anonymitet så är det bra att döpa testpersonerna enligt ett system där personerna är indexerade. Exempelvis kan ``tp\_1'' betyda ``testperson 1''. Om alla namn börjar med något som står för testledarna eller experimentet, så kommer alla personer att hamna efter varandra i programmet. I denna guide har jag därför kallat personen ``ledare\_tp\_1''. Byt ut ``ledare'' mot något som är unikt för er.
}

\newpage

%%%%%%%%%%%%%%%%%%%%%%%%%%%%%%%%%%%%%%%%%%%%%%%%%%%%%%%%%%%%%%
%%% IV. Anslut biomonitor | Steg 1/1
%%%%%%%%%%%%%%%%%%%%%%%%%%%%%%%%%%%%%%%%%%%%%%%%%%%%%%%%%%%%%%

\pagetitle{IV. Anslut biomonitor}{Steg 1/1}

\begin{instructionsenv}
	\instruction{1}{Anslut sladden till ``CH1--4''.}
\end{instructionsenv}

\begin{bild}{biomonitor}{0.6}
	\siffrastreckellips{1}{0.55}{0.7}{1.2}{0.8}{0.63}{0.53}
\end{bild}

\kommentarbox{
	ME6000-T16 har fyra ingångar, som med sladdar förgrenar sig till fyra kanaler för varje ingång. Protokollet definierar namn för kanalerna. Om ni exempelvis sätter sladden i ``CH1--4'', så måste protokollet ha namn för kanalerna i omfånget 1--4.
	Sladdarna på HSC är märkta så att varje förgrening har en siffra. Man kan sätta sladdarna i ``fel'' ingång, bara man behåller inbördes ordning. Det vill säga, sätter man sladden med förgreningar märkta 5--8 i ingången märkt 1--4, så kommer förgrening 5 motsvara kanal 1, förgrening 6 motsvara kanal 2, osv.\par
	Slå på biomonitorn genom att hålla inne knappen märkt med en punkt på framsidan nere till vänster.\par
	Datorns operativsystem kan visa ett meddelande att enheten är ansluten. Stäng ner alla program utom MegaWin, inklusive eventuella antivirusprogram,
}

\newpage

%%%%%%%%%%%%%%%%%%%%%%%%%%%%%%%%%%%%%%%%%%%%%%%%%%%%%%%%%%%%%%
%%% V. Kontrollera signalen | Steg 1/4
%%%%%%%%%%%%%%%%%%%%%%%%%%%%%%%%%%%%%%%%%%%%%%%%%%%%%%%%%%%%%%

\pagetitle{V. Kontrollera signalen}{Steg 1/4}

\begin{instructionsenv}
	\instruction{1}{Klicka på ``Measure''.}
\end{instructionsenv}

\begin{bild}{3}{0.91}
	\siffrastreckellips{1}{0.45}{0.65}{2}{2}{0.5}{0.85}
\end{bild}

\kommentarbox{
Att kontrollera signalen är i princip samma som att göra en mätning, men har här fått ett eget kapitel. \\
Det är bra att kontrollera signalen när försökspersonen är uppvärmd och redo för de riktiga mätningarna. Om något händer med elektroderna eller utrustningen under experimentet kan det vara bra att kontrollera signalen igen.
}
\newpage

%%%%%%%%%%%%%%%%%%%%%%%%%%%%%%%%%%%%%%%%%%%%%%%%%%%%%%%%%%%%%%
%%% V. Kontrollera signalen | Steg 2/4
%%%%%%%%%%%%%%%%%%%%%%%%%%%%%%%%%%%%%%%%%%%%%%%%%%%%%%%%%%%%%%

\pagetitle{V. Kontrollera signalen}{Steg 2/4}

\begin{instructionsenv}
	\instruction{1}{Kontrollera att ``Protocol'' och ``Person'' stämmer.}
	\instruction{2}{Klicka på ``Run protocol''.}
\end{instructionsenv}

\begin{bild}{19}{0.91}
	\siffrastreckellips{1}{0.1}{0.7}{1.5}{1.5}{0.1}{0.92}
	\siffrastreckellips{2}{0.29}{0.7}{2}{0.5}{0.29}{0.945}
\end{bild}

\newpage

%%%%%%%%%%%%%%%%%%%%%%%%%%%%%%%%%%%%%%%%%%%%%%%%%%%%%%%%%%%%%%
%%% V. Kontrollera signalen | Steg 3/4
%%%%%%%%%%%%%%%%%%%%%%%%%%%%%%%%%%%%%%%%%%%%%%%%%%%%%%%%%%%%%%

\pagetitle{V. Kontrollera signalen}{Steg 3/4}

\begin{instructionsenv}
	\instruction{1}{Klicka på den röda cirkeln för att starta mätning.}

\end{instructionsenv}

\begin{bild}{20}{0.91}
	\siffrastreckellips{1}{0.85}{0.7}{0.7}{0.7}{0.975}{0.82}
\end{bild}

\newpage

%%%%%%%%%%%%%%%%%%%%%%%%%%%%%%%%%%%%%%%%%%%%%%%%%%%%%%%%%%%%%%
%%% V. Kontrollera signalen | Steg 4/4
%%%%%%%%%%%%%%%%%%%%%%%%%%%%%%%%%%%%%%%%%%%%%%%%%%%%%%%%%%%%%%

\pagetitle{V. Kontrollera signalen}{Steg 4/4}

\begin{instructionsenv}
	\instruction{1}{Inspektera signalerna i vila.}
	\instruction{2}{Inspektera signalerna under kontraktion.}
	\instruction{3}{Klicka på den grå fyrkanten för att stoppa mätningen.\par \hspace{0.045\linewidth} Vid ``Save measured data?'', välj ``No''.}
\end{instructionsenv}

\begin{bild}{21}{0.91}
	\siffrastreckellips{1}{0.17}{0.45}{1}{1}{0.15}{0.65}
	\siffrastreckellips{1}{0.17}{0.45}{1}{1}{0.15}{0.27}
	
	\siffrastreckellips{2}{0.57}{0.45}{2}{2}{0.265}{0.65}
	\siffrastreckellips{2}{0.57}{0.45}{2}{2}{0.39}{0.28}
	
	\siffrastreckellips{3}{0.85}{0.7}{0.7}{0.7}{0.977}{0.77}
\end{bild}

\kommentarbox{
I vila bör signalen vara tunn. Bilden visar en vilosignal som innehåller lite brus, med god förberedelse av huden och med avslappnad muskel kan signalen vara tunnare.\par
Utför lämpliga rörelser/kontraktioner för att se att aktivitet registreras i varje muskel.  
}
\newpage

%%%%%%%%%%%%%%%%%%%%%%%%%%%%%%%%%%%%%%%%%%%%%%%%%%%%%%%%%%%%%%
%%% VI. Utför mätning | Steg 1/4
%%%%%%%%%%%%%%%%%%%%%%%%%%%%%%%%%%%%%%%%%%%%%%%%%%%%%%%%%%%%%%

\pagetitle{VI. Utför mätning}{Steg 1/4}

\begin{instructionsenv}
	\instruction{1}{Klicka på ``Measure''.}
\end{instructionsenv}

\begin{bild}{3}{0.91}
	\siffrastreckellips{1}{0.45}{0.65}{2}{2}{0.5}{0.85}
\end{bild}

\kommentarbox{
Om du precis kontrollerat signalen så är du kanske redan inne i ``Measure''-delen av programmet och kan hoppa fram till steg 3.
}
\newpage

%%%%%%%%%%%%%%%%%%%%%%%%%%%%%%%%%%%%%%%%%%%%%%%%%%%%%%%%%%%%%%
%%% VI. Utför mätning | Steg 2/4
%%%%%%%%%%%%%%%%%%%%%%%%%%%%%%%%%%%%%%%%%%%%%%%%%%%%%%%%%%%%%%

\pagetitle{V. Utför mätning}{Steg 2/4}

\begin{instructionsenv}
	\instruction{1}{Kontrollera att ``Protocol'' och ``Person'' stämmer.}
	\instruction{2}{Klicka på ``Run protocol''.}
\end{instructionsenv}

\begin{bild}{19}{0.91}
	\siffrastreckellips{1}{0.1}{0.7}{1.5}{1.5}{0.1}{0.92}
	\siffrastreckellips{2}{0.29}{0.7}{2}{0.5}{0.29}{0.945}
\end{bild}

\newpage

%%%%%%%%%%%%%%%%%%%%%%%%%%%%%%%%%%%%%%%%%%%%%%%%%%%%%%%%%%%%%%
%%% VII. Utför mätning | Steg 3/4
%%%%%%%%%%%%%%%%%%%%%%%%%%%%%%%%%%%%%%%%%%%%%%%%%%%%%%%%%%%%%%

\pagetitle{VII. Utför mätning}{Steg 3/4}

\begin{instructionsenv}
	\instruction{1}{Klicka på den röda cirkeln för att starta mätning.}
	
\end{instructionsenv}

\begin{bild}{20}{0.91}
	\siffrastreckellips{1}{0.85}{0.7}{0.7}{0.7}{0.975}{0.82}
\end{bild}

\newpage

%%%%%%%%%%%%%%%%%%%%%%%%%%%%%%%%%%%%%%%%%%%%%%%%%%%%%%%%%%%%%%
%%% VII. Utför mätning | Steg 4/4
%%%%%%%%%%%%%%%%%%%%%%%%%%%%%%%%%%%%%%%%%%%%%%%%%%%%%%%%%%%%%%

\pagetitle{VII. Utför mätning}{Steg 4/4}

\begin{instructionsenv}
	\instruction{1}{Klicka på flaggan för att placera en tidsmarkör.\par  \hspace{0.045\linewidth} Klicka fler gånger för att placera ytterligare tidsmarkörer.}
	\instruction{2}{Klicka på den grå fyrkanten för att stoppa mätningen.\par  \hspace{0.045\linewidth} Om du är nöjd med mätningen, spara denna data.}
	
\end{instructionsenv}

\begin{bild}{23}{0.91}
	\siffrastreckellips{1}{0.85}{0.53}{0.6}{0.7}{0.975}{0.53}
	
	\node[black, thick] (A) at  (0.26,0.75) {Tidsmarkörer};
	\node (B) at  (0.34,0.83) {};
	\node (C) at  (0.19,0.83) {};
	\draw[-Stealth, black, line width=1pt] (A)--(B);
	\draw[-Stealth, black, line width=1pt] (A)--(C);
	
	\siffrastreckellips{2}{0.85}{0.7}{0.7}{0.7}{0.977}{0.77}
\end{bild}

\kommentarbox{
	I detta exempel har jag satt två flaggor för att markera start- och  stopptiden då försökspersonen utförde en viss rörelse. Om man vill kan man låta mätningen fortsätta och sätta fler flaggor vid nästa rörelse. Det gäller då att man är strukturerad så att man senare vet vad som hände mellan flagga 1 och 2, och mellan flagga 3 och 4, osv. För enkelhetens skull nöjer jag mig med två flaggor här.\par 
	Du har nu utfört en mätning. Om du tänkt göra fler mätningar innan du analyserar datan, så är det en god idé att ge mätningen en kommentar som identifierar vad du precis mätt. Hur man gör detta kan du läsa i nästa kapitel.
}

\newpage

%%%%%%%%%%%%%%%%%%%%%%%%%%%%%%%%%%%%%%%%%%%%%%%%%%%%%%%%%%%%%%
%%% VII. Bearbeta signalen | Steg 1/17
%%%%%%%%%%%%%%%%%%%%%%%%%%%%%%%%%%%%%%%%%%%%%%%%%%%%%%%%%%%%%%

\pagetitle{VII. Bearbeta signalen}{Steg 1/17}

\begin{instructionsenv}
	\instruction{1}{Klicka på ``Results''.}
\end{instructionsenv}

\begin{bild}{3}{0.91}
	\siffrastreckellips{1}{0.8}{0.65}{2}{2}{0.863}{0.85}
\end{bild}

\kommentarbox{
	Denna sektion kommer att visa ett exempel på hur data kan bearbetas. Låt oss säga att test\-ledaren har instruerat testpersonen att utföra en isometrisk kontraktion under fem sekunder, och satt tidsmarkörer vid början och slutet av de fem sekunderna. Vi tänker oss vidare att vi ska använda data från tre sekunder i mitten av kontraktionen. Från denna data vill vi se medelaktivitet i musklerna.
}

\newpage

%%%%%%%%%%%%%%%%%%%%%%%%%%%%%%%%%%%%%%%%%%%%%%%%%%%%%%%%%%%%%%
%%% VII. Bearbeta signalen | Steg 2/17
%%%%%%%%%%%%%%%%%%%%%%%%%%%%%%%%%%%%%%%%%%%%%%%%%%%%%%%%%%%%%%

\pagetitle{VII. Bearbeta signalen}{Steg 2/17}

\begin{instructionsenv}
	\instruction{1}{Högerklicka på fältet under ``Comment'' och ange en kommentar.}
	\instruction{2}{Dubbelklicka på ``Raw free''}
\end{instructionsenv}

\begin{bild}{25}{0.91}
	\siffrastreckellips{1}{0.34}{0.75}{2}{0.3}{0.34}{0.87}
	\siffrastreckellips{2}{0.81}{0.75}{3}{0.4}{0.81}{0.897}
\end{bild}

\kommentarbox{
	I detta fiktiva experiment tänkte jag mig att vi precis mätt första försöket av övning ett, och kallade därför denna mätning för ``Övning 1, försök 1''.\par 
	``Raw free är den obearbetade signalen ni sparade från mätningen. I följande steg kommer vi att bearbeta signalen med tidsbegränsning, filter, rektifiering, mm. Ett antal olika snarlika fönster kommer att öppnas, och det gäller att hålla koll så att man gör rätt sak på rätt plats, och sparar rätt data.
}

\newpage

%%%%%%%%%%%%%%%%%%%%%%%%%%%%%%%%%%%%%%%%%%%%%%%%%%%%%%%%%%%%%%
%%% VII. Bearbeta signalen | Steg 3/17
%%%%%%%%%%%%%%%%%%%%%%%%%%%%%%%%%%%%%%%%%%%%%%%%%%%%%%%%%%%%%%

\pagetitle{VII. Bearbeta signalen}{Steg 3/17}

\begin{instructionsenv}
	\instruction{1}{Klicka ``Zoom from marker to marker''.}
\end{instructionsenv}

\vspace{-20pt}

\begin{bild}{26}{0.91}
	\siffrastreckellips{1}{0.272}{1.08}{0.7}{0.7}{0.272}{0.894}
\end{bild}

\kommentarbox{
	Här ska vi med hjälp av flaggorna välja ut vilket tidsomfång som vi är intresserade av.
}

\newpage

%%%%%%%%%%%%%%%%%%%%%%%%%%%%%%%%%%%%%%%%%%%%%%%%%%%%%%%%%%%%%%
%%% VII. Bearbeta signalen | Steg 4/17
%%%%%%%%%%%%%%%%%%%%%%%%%%%%%%%%%%%%%%%%%%%%%%%%%%%%%%%%%%%%%%

\pagetitle{VII. Bearbeta signalen}{Steg 4/17}

\begin{instructionsenv}
	\instruction{1}{Markera två flaggor.}
	\instruction{2}{Klicka ``Ok''.}
\end{instructionsenv}

\begin{bild}{27}{0.91}
	\siffrastreckellips{1}{0.4}{0.5}{0.9}{0.9}{0.305}{0.58}
	\siffrastreckellips{2}{0.55}{0.37}{1.4}{0.5}{0.65}{0.22}	
\end{bild}

\kommentarbox{
	Om ni har satt fler än två flaggor under samma mätning gäller det att ni nu minns vilka flaggor ni är intresserade av.
}

\newpage

%%%%%%%%%%%%%%%%%%%%%%%%%%%%%%%%%%%%%%%%%%%%%%%%%%%%%%%%%%%%%%
%%% VII. Bearbeta signalen | Steg 5/17
%%%%%%%%%%%%%%%%%%%%%%%%%%%%%%%%%%%%%%%%%%%%%%%%%%%%%%%%%%%%%%

\pagetitle{VII. Bearbeta signalen}{Steg 5/17}

\begin{instructionsenv}
	\instruction{1}{Klicka ``Edit data''.}
\end{instructionsenv}

\vspace{-20pt}

\begin{bild}{28}{0.91}
	\siffrastreckellips{1}{0.502}{1.08}{0.7}{0.7}{0.502}{0.894}
\end{bild}

\kommentarbox{
	Fönstret är nu fyllt av aktiviteten från flagga 1 till flagga 2. Tiden mellan flaggorna är ungefär 5 sekunder, vilket kan avläsas på tidsskalan nedtill. De streckade horisontella linjerna markerar sekunder.
}

\newpage

%%%%%%%%%%%%%%%%%%%%%%%%%%%%%%%%%%%%%%%%%%%%%%%%%%%%%%%%%%%%%%
%%% VII. Bearbeta signalen | Steg 6/17
%%%%%%%%%%%%%%%%%%%%%%%%%%%%%%%%%%%%%%%%%%%%%%%%%%%%%%%%%%%%%%

\pagetitle{VII. Bearbeta signalen}{Steg 6/17}

\begin{instructionsenv}
	\instruction{1}{Klicka ``Digital filtering''.}
\end{instructionsenv}

\vspace{-20pt}

\begin{bild}{29}{0.91}
	\siffrastreckellips{1}{0.582}{1.08}{0.7}{0.7}{0.582}{0.923}
\end{bild}

\kommentarbox{
	Vi kan börja med att applicera frekvensfilter. Här kommer jag att visa hur man filtrerar med bandpass och notch-filter.
}

\newpage

%%%%%%%%%%%%%%%%%%%%%%%%%%%%%%%%%%%%%%%%%%%%%%%%%%%%%%%%%%%%%%
%%% VII. Bearbeta signalen | Steg 7/17
%%%%%%%%%%%%%%%%%%%%%%%%%%%%%%%%%%%%%%%%%%%%%%%%%%%%%%%%%%%%%%

\pagetitle{VII. Bearbeta signalen}{Steg 7/17}

\begin{instructionsenv}
	\instruction{1}{Välj ``Band pass''.}
	\instruction{2}{Fyll i ``Low frequency'': 15 Hz.}
	\instruction{3}{Fyll i ``High frequency'': 400 Hz.}
	\instruction{4}{Klicka ``OK''.}		
\end{instructionsenv}

\begin{bild}{30}{0.91}
	\siffrastreckellips{1}{0.6}{0.88}{2}{0.4}{0.54}{0.746}
	\siffrastreckellips{2}{0.67}{0.635}{1.7}{0.4}{0.53}{0.635}
	\siffrastreckellips{3}{0.67}{0.57}{1.7}{0.4}{0.53}{0.57}
	\siffrastreckellips{4}{0.945}{0.18}{1.5}{0.6}{0.945}{0.032}		
\end{bild}

\kommentarbox{
	Var noga med att välja ``Band pass'' och inte ``Band stop''.
}

\newpage

%%%%%%%%%%%%%%%%%%%%%%%%%%%%%%%%%%%%%%%%%%%%%%%%%%%%%%%%%%%%%%
%%% VII. Bearbeta signalen | Steg 8/17
%%%%%%%%%%%%%%%%%%%%%%%%%%%%%%%%%%%%%%%%%%%%%%%%%%%%%%%%%%%%%%

\pagetitle{VII. Bearbeta signalen}{Steg 8/17}

\begin{instructionsenv}
	\instruction{1}{Klicka ``Digital filtering'' igen.} 
\end{instructionsenv}

\vspace{-20pt}

\begin{bild}{29}{0.91}
	\siffrastreckellips{1}{0.582}{1.08}{0.7}{0.7}{0.582}{0.923}
\end{bild}

\kommentarbox{
	Vi hittade inget sätt att applicera båda filtren på en gång.
}

\newpage

%%%%%%%%%%%%%%%%%%%%%%%%%%%%%%%%%%%%%%%%%%%%%%%%%%%%%%%%%%%%%%
%%% VII. Bearbeta signalen | Steg 9/17
%%%%%%%%%%%%%%%%%%%%%%%%%%%%%%%%%%%%%%%%%%%%%%%%%%%%%%%%%%%%%%

\pagetitle{VII. Bearbeta signalen}{Steg 9/17}

\begin{instructionsenv}
	\instruction{1}{Välj ``Notch 50Hz''.}
	\instruction{2}{Klicka ``OK''.}		
\end{instructionsenv}

\begin{bild}{32}{0.91}
	\siffrastreckellips{1}{0.6}{0.88}{2}{0.4}{0.54}{0.746}
	\siffrastreckellips{2}{0.945}{0.18}{1.5}{0.6}{0.945}{0.032}		
\end{bild}

\kommentarbox{
	Standardspänningen i alla vanliga enfas vägguttag i Sverige är en växelspänning på 50 Hz.
}

\newpage

%%%%%%%%%%%%%%%%%%%%%%%%%%%%%%%%%%%%%%%%%%%%%%%%%%%%%%%%%%%%%%
%%% VII. Bearbeta signalen | Steg 10/17
%%%%%%%%%%%%%%%%%%%%%%%%%%%%%%%%%%%%%%%%%%%%%%%%%%%%%%%%%%%%%%

\pagetitle{VII. Bearbeta signalen}{Steg 10/17}

\begin{instructionsenv}
	\instruction{1}{Stäng fönstret. Vid ``Save changes?'', välj ``Yes''.}
\end{instructionsenv}

\vspace{-20pt}

\begin{bild}{34}{0.91}
	\siffrastreckellips{1}{0.984}{1.08}{0.4}{0.4}{0.984}{0.98}
\end{bild}

\newpage

%%%%%%%%%%%%%%%%%%%%%%%%%%%%%%%%%%%%%%%%%%%%%%%%%%%%%%%%%%%%%%
%%% VII. Bearbeta signalen | Steg 11/17
%%%%%%%%%%%%%%%%%%%%%%%%%%%%%%%%%%%%%%%%%%%%%%%%%%%%%%%%%%%%%%

\pagetitle{VII. Bearbeta signalen}{Steg 11/17}

\begin{instructionsenv}
	\hspace{0.045\linewidth} \parbox{0.955\linewidth}{{\large Nu har ni förmodligen en hel del fönster öppna med data i olika stadier av bearbetning. Stäng ner dem tills ni ser fönstret nedan. Ni behöver inte spara datan från dessa fönster.}}\par 
	\instruction{1}{Dubbelklicka på ``Edited data''.}
\end{instructionsenv}

\begin{bild}{35}{0.91}
	\siffrastreckellips{1}{0.81}{0.75}{3}{0.4}{0.81}{0.87}
\end{bild}

\kommentarbox{
	``Edited data'' är den filtrerade datan från föregående steg. 
}

\newpage

%%%%%%%%%%%%%%%%%%%%%%%%%%%%%%%%%%%%%%%%%%%%%%%%%%%%%%%%%%%%%%
%%% VII. Bearbeta signalen | Steg 12/17
%%%%%%%%%%%%%%%%%%%%%%%%%%%%%%%%%%%%%%%%%%%%%%%%%%%%%%%%%%%%%%

\pagetitle{VII. Bearbeta signalen}{Steg 12/17}

\begin{instructionsenv} 
	\instruction{1}{Klicka på på ``RMS Averaging''.}
\end{instructionsenv}

\begin{bild}{36}{0.91}
	\siffrastreckellips{1}{0.2}{0.7}{0.7}{0.7}{0.22}{0.83}
\end{bild}

\kommentarbox{
	RMS är ett sätt att bearbeta datan för att få positiva värden som kan användas för uträkningar av medelvärden.\par 
	RMS står för Root Mean Squared, och för den matematiskt intresserade läsaren kanske namnet påminner om det statistiska begreppet standardavvikelse; RMS är helt enkelt samma sak som standardavvikelse i fallet då summan av medelvärdena är noll. RMS beräknas med ett glidande fönster.
}

\newpage

%%%%%%%%%%%%%%%%%%%%%%%%%%%%%%%%%%%%%%%%%%%%%%%%%%%%%%%%%%%%%%
%%% VII. Bearbeta signalen | Steg 13/17
%%%%%%%%%%%%%%%%%%%%%%%%%%%%%%%%%%%%%%%%%%%%%%%%%%%%%%%%%%%%%%

\pagetitle{VII. Bearbeta signalen}{Steg 13/17}

\begin{instructionsenv}
	\hspace{0.045\linewidth} \parbox{0.955\linewidth}{{\large Vi bestämde tidigare att detta experiment ska använda de tre mellersta sekund\-erna som grund för resultaten.}}\par
	\instruction{1}{Klicka på ``AOI''.}
	\instruction{2}{Klicka och dra muspekaren över ett tidsomfång.}
\end{instructionsenv}

\begin{bild}{37}{0.91}
	\siffrastreckellips{1}{0.317}{0.7}{0.7}{0.7}{0.317}{0.891}
	\node[mycolor1] (A) at (0.252,0.58){};
	\node[mycolor1, thick] (B) at (0.541,0.58){\LARGE 2};
	\node[mycolor1] (C) at (0.83,0.58){};
	\draw[|-,mycolor1, line width=1pt] (A)--(B);
	\draw[-Stealth,mycolor1, line width=1pt] (B)--(C);
	
	
\end{bild}

\kommentarbox{
``AOI'' står för Area of interest.
}

\newpage

%%%%%%%%%%%%%%%%%%%%%%%%%%%%%%%%%%%%%%%%%%%%%%%%%%%%%%%%%%%%%%
%%% VII. Bearbeta signalen | Steg 14/17
%%%%%%%%%%%%%%%%%%%%%%%%%%%%%%%%%%%%%%%%%%%%%%%%%%%%%%%%%%%%%%

\pagetitle{VII. Bearbeta signalen}{Steg 14/17}

\begin{instructionsenv}
	\instruction{1}{Klicka på ``Basic results''.}
\end{instructionsenv}

\begin{bild}{38}{0.91}
	\siffrastreckellips{1}{0.16}{0.67}{0.8}{0.8}{0.035}{0.83}
\end{bild}

\newpage

%%%%%%%%%%%%%%%%%%%%%%%%%%%%%%%%%%%%%%%%%%%%%%%%%%%%%%%%%%%%%%
%%% VII. Bearbeta signalen | Steg 15/17
%%%%%%%%%%%%%%%%%%%%%%%%%%%%%%%%%%%%%%%%%%%%%%%%%%%%%%%%%%%%%%

\pagetitle{VII. Bearbeta signalen}{Steg 15/17}

\begin{instructionsenv}
	\hspace{0.045\linewidth} \parbox{0.955\linewidth}{{\large Här är resultaten.\\ Om ni vill spara resultaten i en textfil kan ni göra det på följande sätt: }}\par
	\instruction{1}{Klicka på ``ASCII file''.}
\end{instructionsenv}

\begin{bild}{39}{0.91}
	\siffrastreckellips{1}{0.565}{0.68}{0.8}{0.8}{0.565}{0.872}
\end{bild}

\newpage

%%%%%%%%%%%%%%%%%%%%%%%%%%%%%%%%%%%%%%%%%%%%%%%%%%%%%%%%%%%%%%
%%% VII. Bearbeta signalen | Steg 16/17
%%%%%%%%%%%%%%%%%%%%%%%%%%%%%%%%%%%%%%%%%%%%%%%%%%%%%%%%%%%%%%

\pagetitle{VII. Bearbeta signalen}{Steg 16/17}

\begin{instructionsenv}
	\instruction{1}{Klicka på ``OK''.}
\end{instructionsenv}

\begin{bild}{40}{0.91}
	\siffrastreckellips{1}{0.5}{0.3}{1.7}{0.7}{0.6}{0.175}
\end{bild}

\newpage

%%%%%%%%%%%%%%%%%%%%%%%%%%%%%%%%%%%%%%%%%%%%%%%%%%%%%%%%%%%%%%
%%% VII. Bearbeta signalen | Steg 17/17
%%%%%%%%%%%%%%%%%%%%%%%%%%%%%%%%%%%%%%%%%%%%%%%%%%%%%%%%%%%%%%

\pagetitle{VII. Bearbeta signalen}{Steg 17/17}

\begin{instructionsenv}
	\instruction{1}{Ange ett lämpligt namn för filen.}
	\instruction{2}{Klicka på ``Save''.}
\end{instructionsenv}

\begin{bild}{41}{0.91}
	\siffrastreckellips{1}{0.3}{0.4}{4}{0.9}{0.32}{0.2}
	\siffrastreckellips{2}{0.7}{0.4}{2.8}{0.8}{0.875}{0.2}
\end{bild}

\kommentarbox{
	Välj gärna ett filnamn som är unikt för testpersonen och mätningen, i en mapp som ni sedan kan kopiera och ta med.
}

\newpage

\end{document}